% Supplementary literature positioning for the SIVP manuscript.
\appendix
\section{Additional Literature Positioning}
\label{app:literature-positioning}

This appendix briefly positions RA-RepDet within broader detection, multimodal fusion, event perception, evaluation-integrity, and lightweight-model literature. It is included as supporting context and does not expand the experimental claims beyond the component-disjoint development-validation and locked same-dataset held-out guard evidence reported in the main text.

Generic object-detection and aerial-detection studies motivate the use of multi-scale lightweight detectors for UAV scenes with small objects, changing viewpoints, clutter, and scale variation. Visible-thermal and event-based perception studies motivate RGB--thermal--event fusion, but the present work uses an already constructed five-channel representation rather than claiming raw event-stream reconstruction or sensor-synchronization optimization.

Missing-modality and multimodal-regularization work provides context for modality dropout. In this manuscript, modality dropout is a training-time intervention and any channel-zeroing result should be interpreted only as a controlled stress test. It is not a physical sensor-failure model and does not establish calibrated sensor reliability.

Evaluation-leakage and dataset-quality studies motivate the component-disjoint protocol. The held-out guard evaluation reduces reliance on a single development-validation split, but it remains within the TriAir project dataset and should not be described as external-dataset generalization or as an independent public benchmark test.
